\documentclass[10pt]{article}
% Math Packages
\usepackage{amsmath, mathtools}
\usepackage{amssymb}
\usepackage{amsthm}
\usepackage{amsfonts}
\usepackage{bbm}
\usepackage{breqn}
\usepackage[margin=1in]{geometry}
\usepackage{graphicx}
\usepackage{tikz}
\usepackage{forest}
\usepackage{tikz-qtree}
\graphicspath{ {./images/} }
\usepackage{hyperref}
\usepackage[capitalize]{cleveref}
\usepackage[shortlabels]{enumitem}
\usetikzlibrary{arrows,matrix,positioning}

% Colorful Notes
\usepackage{color} \definecolor{Red}{rgb}{1,0,0} \definecolor{Blue}{rgb}{0,0,1}
\definecolor{Purple}{rgb}{.5,0,.5} \def\red{\color{Red}} \def\blue{\color{Blue}}
\def\gray{\color{gray}} \def\purple{\color{Purple}}
\newcommand{\rnote}[1]{{\red [#1]}} % \rnote{foo} gives '[foo]' in red
\newcommand{\pnote}[1]{{\purple [#1]}} % \pnote{foo} gives '[foo]' in purple
\newcommand{\bnote}[1]{{\blue #1}} % \bnote{foo} gives 'foo' in blue
\newcommand{\gnote}[1]{{\gray #1}} % \gnote{foo} gives 'foo' in gray
\newcommand{\Max}[1]{{\purple [#1]}} % \bnote{foo} then 'foo' is blue
\definecolor{RoyalBlue}{cmyk}{1, 0.50, 0, 0}
\newcommand{\dempfcolor}[1]{{\color{RoyalBlue}#1}} 
\newcommand{\demph}[1]{\dempfcolor{{\sl #1}}}




% Math Environments
\newtheorem{theorem}{Theorem}
\newtheorem{definition}[theorem]{Definition}
\newtheorem{assumption}[theorem]{Assumption}
\newtheorem{lemma}[theorem]{Lemma}
\newtheorem{remark}[theorem]{Remark}
\newtheorem{proposition}[theorem]{Proposition}
\newtheorem{corollary}[theorem]{Corollary}
\newtheorem{example}[theorem]{Example} 
\newtheorem{question}[theorem]{Question}

% Custom Math Commands
\newcommand{\vt}{\vskip 5mm} % vertical space
\newcommand{\fl}{\\\noindent\textbf} % first line
\newcommand{\Fl}{\vt\noindent\textbf} % first line with space above


\begin{document}
\section*{Syllabus for Math 243: Calculus III (Section 001)}
Spring 2025 (Jan 13 - May 16)

\Fl{Instructor:} Max Hill
\fl{Office:} Physical Sciences Building 304
\fl{Office Hours:} MWF 2-3pm (or other times by appointment).

\Fl{Place and time:} Watanabe 112 at 10:30-11:20am (MWF)

\Fl{Textbook:} \textit{Calculus} by James Stewart (8th edition). ISBN-13: \texttt{978-1-305-26665-0}

\Fl{Course prerequistes:} a grade of C or better in Math 242 or Math 252A. If you don't meet the
prerequisites, you'll need to get approval from the math department office.


\Fl{Grading:} Homework (30\%), Midterm (35\%), Final (35\%). The following grade cutoffs will be used at
the end of the semester to determine final grades:
\begin{center}
  \begin{tabular}{|c|c|c|c|c|c|c|c|c|c|c|}
    \hline
    D&D+&C-&C&C+&B-&B&B+&A-&A&A+  \\
    \hline
    60\%&67\%&70\%&73\%&77\%&80\%&83\%&87\%&90\%&93\%&97\%\\
    \hline
  \end{tabular}
\end{center}
I may change the above numbers, but if I do so it will be in your favor.



\Fl{Homework:} Homework will be due approximately weekly. 
\begin{itemize}
  \item You have one free `no questions asked' homework extension.
  \item Your lowest homework score will be dropped at the end of the semester.
  \item You may collaborate with classmates on the homeworks. But if you do so, you must (1) make an
  effort write up your solutions on your own, using your own words, and (2) list the names of the people
  who you worked with.
\end{itemize}

\Fl{Exams:} There will be a midterm and a final exam. The final will be cumulative, but will emphasize
material after the midterm. The exam dates are as follows:
\begin{center} \begin{minipage}{5.0in}
    \begin{flushleft}
      Midterm exam   \dotfill~Monday, March 24 (in class) \\
      Final exam        \dotfill~Monday, May 12, 9:45-11:45am\\
    \end{flushleft}
  \end{minipage}
\end{center}
If you do better on the final than the midterm, then your grade on the final will replace your midterm
grade.

\Fl{Make-up policy:} Make-up exams are allowed only in three types of circumstances: (1) in accordance with
university policies, such as conflict with a religious observation, (2) conflicts with another
university-related event, or (3) exceptional circumstances, such as a last-minute medical or family
emergency with verification. In the first two cases, notice must be given to the instructor two weeks in
advance.

\Fl{Accommodations:} Any student who feels s/he may need an accommodation based on the impact of a
disability is invited to contact me privately. I would be happy to work with you, and the KOKUA Program
(Office for Students with Disabilities) to ensure reasonable accommodations in my course. KOKUA can be
reached at (808) 956-7511 or (808) 956-7612 (voice/text) in Room 013 of the Queen Lili`uokalani Center
for Student Services.




\newpage
\Fl{Tentative course outline:}

\begin{itemize}
  \item \textbf{Weeks 1-3:} Parametric equations and Polar coordinates
  \begin{enumerate}
    \item Section 8.1: Arc length
    \item  Section 10.1: Curves defined by parametric equations.
    \item Section 10.2: Calculus with plane curves.
    \item  Section 10.3 Polar coordinates
    \item  Section 10.4: Areas and lengths in polar coordinates
    \item Section 10.5: Conic sections.
  \end{enumerate}
  \item \textbf{Weeks 4-6:} Vectors and the geometry of space.
  \begin{enumerate}
    \item Section 12.1: Three-dimensional coordinate systems.
    \item Section 12.2: Vectors.
    \item  Section 12.3: The dot product.
    \item  Section 12.4: The cross product.
    \item Section 12.5: Equations of lines and planes.
    \item  Section 12.6: Cylinders and quadric surfaces.
  \end{enumerate}
  \item \textbf{Weeks 7-9:} Vector functions.
  \begin{enumerate}
    \item Section 13.1: Vector functions and space curves
    \item Section 13.2: Derivatives and integrals of vector functions.
    \item Section 13.3: Arc length and curvature.
    \item Section 13.4: Motion in space: velocity and acceleration.
  \end{enumerate}
  \item \textbf{Weeks 10-15:} Partial derivatives.
  \begin{enumerate}
    \item Section 14.1: Functions of several variables.
    \item Section 14.2: Limits and continuity
    \item Section 14.3: Partial derivatives.
    \item Section 14.4 Tangent planes and linear approximations
    \item Section 14.5: The chain rule.
    \item Section 14.6: Directional derivatives and the gradient vector.
    \item Section 14.7: Maximum and minimum values.
    \item Discovery project p.1010: Quadratic approximation and critical
    points (Taylor’s formula for two variables).
    \item Section 14.8: Lagrange multipliers.
  \end{enumerate}
\end{itemize} 
\end{document}